\lysh{20655_SOLUTION}{1}
\hfill \break
\textbf{α)}
\begin{enumerate}[i.]
    \item Για να αποδείξουμε ότι τα σημεία Α, Β και Γ δεν είναι συνευθειακά αρκεί να δείξουμε ότι $\vec{AB} \nmid \vec{A\Gamma}$.

    Είναι
    \[
        \vec{AB} = (x_B - x_A, y_B - y_A) = (3-2, -1-1) = (1,-2)
    \]
    \[
        \vec{A\Gamma} = (x_\Gamma - x_A, y_\Gamma - y_A) = (-2-2, 0-1) = (-4,-1)
    \]
    και
    \[
        \det(\vec{AB}, \vec{A\Gamma}) = \begin{vmatrix} 1 & -2 \\ -4 & -1 \end{vmatrix} = 1 \cdot (-1) - (-2) \cdot (-4) = -1 - 8 = -9
    \]
    Επειδή $\det(\vec{AB}, \vec{A\Gamma}) \ne 0$ είναι $\vec{AB} \nmid \vec{A\Gamma}$ και έπεται το ζητούμενο.

    \item Είναι $\left( \text{AB}\Gamma \right) = \frac{1}{2} \left| \det(\vec{AB}, \vec{A\Gamma}) \right| = \frac{1}{2} \left| -9 \right| = \frac{9}{2}$ τ.μ.
\end{enumerate}

\textbf{β)} Το εμβαδό του τριγώνου $\Delta \text{A}\Gamma$ ισούται με
\[
    (\Delta \text{A}\Gamma) = \frac{1}{2} |\det(\vec{A\Delta}, \vec{A\Gamma})| = \frac{1}{2} \left| \begin{vmatrix} x-2 & y-1 \\ -4 & -1 \end{vmatrix} \right|
\]
\[
    = \frac{1}{2} |(x-2)(-1) - (y-1)(-4)|
\]
\[
    = \frac{1}{2} |-x+2+4y-4| = \frac{1}{2} |-x+4y-2|
\]
Επομένως, το $\Delta(x,y)$ είναι σημείο του γεωμετρικού τόπου, αν και μόνο αν ισχύει
\[
    (\Delta \text{A}\Gamma) = (\text{AB}\Gamma) \iff \frac{1}{2} |-x+4y-2| = \frac{9}{2}
\]
\[
    \iff |-x+4y-2| = 9
\]
\[
    \iff -x+4y-2 = 9 \text{ ή } -x+4y-2 = -9
\]
\[
    \iff -x+4y-11 = 0 \text{ ή } -x+4y+7 = 0
\]
Άρα, ο ζητούμενος γεωμετρικός τόπος αποτελείται από τις ευθείες $x-4y-7=0$ και $x-4y+11=0$.

\textbf{γ)}
\begin{enumerate}[i.]
    \item Είναι $\lambda_{\text{A}\Gamma} = \frac{y_\Gamma - y_\text{A}}{x_\Gamma - x_\text{A}} = \frac{0-1}{-2-2} = \frac{-1}{-4} = \frac{1}{4}$ και $\lambda_{\epsilon_1} = - \frac{1}{-4} = \frac{1}{4}$. Επειδή οι ευθείες $\epsilon_1, \epsilon_2$ και $\text{A}\Gamma$ έχουν ίσους συντελεστές διεύθυνσης είναι μεταξύ τους παράλληλες.

    \item \textbf{α τρόπος}

    Είναι
    \[
        d(\Gamma, \epsilon_1) = \frac{|x_\Gamma - 4y_\Gamma - 7|}{\sqrt{1^2+(-4)^2}} = \frac{|-2 - 4 \cdot 0 - 7|}{\sqrt{1+16}} = \frac{|-9|}{\sqrt{17}} = \frac{9\sqrt{17}}{17} \text{ μονάδες μήκους}
    \]
    \[
        d(\Gamma, \epsilon_2) = \frac{|x_\Gamma - 4y_\Gamma + 11|}{\sqrt{1^2+(-4)^2}} = \frac{|-2 - 4 \cdot 0 + 11|}{\sqrt{1+16}} = \frac{|9|}{\sqrt{17}} = \frac{9\sqrt{17}}{17} \text{ μονάδες μήκους}
    \]
    και επειδή οι ευθείες $\epsilon_1, \epsilon_2$ και $\text{A}\Gamma$ είναι μεταξύ τους παράλληλες συμπεραίνουμε ότι ο ισχυρισμός «οι ευθείες $x-4y-7=0$ και $x-4y+11=0$ έχουν ως μεσοπαράλληλο την ευθεία $\text{A}\Gamma$» είναι αληθής.

    \textbf{β τρόπος}

    Από το ερώτημα (β) γνωρίζουμε ότι ισχύει $(\Delta \text{A}\Gamma) = (\text{AB}\Gamma)$ μόνο όταν το $\Delta$ ανήκει στην ευθεία $\epsilon_1$ ή στην ευθεία $\epsilon_2$. Αν επιλέξουμε το $\Delta$ να ανήκει στην ευθεία $\epsilon_1$, τότε έχουμε:
    \[
        (\Delta \text{A}\Gamma) = \frac{1}{2} d(\Delta, \text{A}\Gamma) \cdot |\vec{A\Gamma}| = \frac{1}{2} d(B, \text{A}\Gamma) \cdot |\vec{A\Gamma}|
    \]
    \[
        \implies d(\Delta, \text{A}\Gamma) = d(B, \text{A}\Gamma)
    \]
    \[
        \implies d(\epsilon_1, \text{A}\Gamma) = d(B, \text{A}\Gamma)
    \]
    Οπότε, ο ισχυρισμός «οι ευθείες $x-4y-16=0$ και $x-4y+20=0$ έχουν ως μεσοπαράλληλο την ευθεία $\text{A}\Gamma$» είναι αληθής.
\end{enumerate}

% TODO: TikZ σχήμα (μην το κατασκευάσεις)
\end{lysh}